\section{专门机关}
\begin{enumerate}
  \item 诉讼参与人一定要把这个\zd{公检法的工作人员}踢出去，他们都属于专门机关。
  \item 权力配置。
        \begin{enumerate}
          \item 公安。双重属性、4 层组织体系、领导关系。职权范围。
                \begin{enumerate}
                  \item 立案权。
                  \item 侦查权。9 种手段，5 种措施。
                  \item 执行权。
                        \begin{enumerate}
                          \item 刑罚的执行。余刑\zd{不足 3 月}、\zd{拘}役、剥\zd{夺政}治权利。
                          \item 5 种强制措施。拘传、取保候审、监视居住、拘留、逮捕。
                        \end{enumerate}
                  \item 公安一般不和法院直接打交道。
                \end{enumerate}
          \item 监狱。是主要的刑罚的执行机关。
                \begin{enumerate}
                  \item \zd{徒刑监禁 3 月起}。\al{管制 1 年 \cuo}。
                  \item 侦查权 \yjt \zd{罪犯在监狱内的犯罪}。\al{监狱监管人员指使体罚虐待 \yjt 职务犯罪}！！
                  \item \zd{提}出监外执行、减刑、假释的\zd{书面意见/建议}。
                \end{enumerate}
          \item 检察院。4 层组织体系、领导关系。\zd{上命下从}。\al{最高检可决定检察院在二审中如何发表意见 \dui}。
                \begin{enumerate}
                  \item 军事检察院与地方检察院平级。
                  \item 检察长对检察员。要求复核 \hz 直接作出决定 \hz 提交检委会来讨论决定。
                  \item 检察员对检察长。书面提出意见。
                  \item \zd{捕诉合一}（批捕、起诉）的改革。原则是同一个检察官/办案组。例外。改变管辖 \hz 另行指派。
                        \al{由于受害者大多在 B 市，A 市检可将案件移送给 B 市检}。
                  \item 职权范围。\zd{立案、侦查权、公诉权、法律监督权}（唯一）。
                  \item 先这么记吧。\zd{除了抗诉不当要先听意见，其他都是可以、直接、能够、决定}。
                \end{enumerate}
          \item 法院。4 层组织体系、单个独立、监督指导关系。\zx{上可审下，下不可审上}。\al{上级法院认为下级审更合适，将案件交由下级法院审理 \cuo}。
                \begin{enumerate}
                  \item 监督关系。除了\zd{改变管辖、二审/再审撤销不合法判决}。其他统统是干涉。\al{讨论、提出意见 \cuo}。
                  \item 不告不理。法院一般永远不主动。
                  \item 执行权。仅限财产罚。
                \end{enumerate}
        \end{enumerate}
\end{enumerate}
